\title{2022 年高考数学 新高考II卷}
\date{}
\maketitle

\section*{一、单项选择题}
本题共8小题，每小题5分，共40分．在每小题给出的四个选项中，只有一项是符合题目要求的.

\begin{question}[points=5]
已知集合 $A=\{-1,1,2,4\}, B=\{x||x-1|\leq 1\}$，则 $A\cap B=$（ ）
\begin{choices}
  \item $\{-1,2\}$
  \item $\{1,2\}$
  \item $\{1,4\}$
  \item $\{-1,4\}$
\end{choices}
\begin{solution}
\textbf{答案：}$\{1,2\}$

\textbf{分析：}求出集合 $B$ 后可求 $A\cap B$.

\textbf{详解：}$B=\{x|0\leq x\leq 2\}$，故 $A\cap B=\{1,2\}$，故选B.
\end{solution}
\end{question}

\begin{question}[points=5]
$(2+2i)(1-2i)=$（ ）
\begin{choices}
  \item $-2+4i$
  \item $-2-4i$
  \item $6+2i$
  \item $6-2i$
\end{choices}
\begin{solution}
\textbf{答案：}$6-2i$

\textbf{分析：}利用复数的乘法可求 $(2+2i)(1-2i)$.

\textbf{详解：}$(2+2i)(1-2i)=2+4-4i+2i=6-2i$，故选D.
\end{solution}
\end{question}

\begin{question}[points=5]
图1是中国古代建筑中的举架结构，$AA',BB',CC',DD'$是桁，相邻桁的水平距离称为步，垂直距离称为举，图2是某古代建筑屋顶截面的示意图．其中$DD_1,CC_1,BB_1,AA_1$是举，$OD_1,DC_1,CB_1,BA_1$是相等的步，相邻桁的举步之比分别为$\dfrac{DD_1}{OD_1}=0.5,\dfrac{CC_1}{DC_1}=k_1,\dfrac{BB_1}{CB_1}=k_2,\dfrac{AA_1}{BA_1}=k_3$．已知$k_1,k_2,k_3$成公差为$0.1$的等差数列，且直线$OA$的斜率为$0.725$，则$k_3=$（ ）

\begin{center}
  \includegraphics[width=.38\linewidth]{question-3-1.jpg}
\end{center}
\begin{choices}
  \item $0.75$
  \item $0.8$
  \item $0.85$
  \item $0.9$
\end{choices}
\begin{solution}
\textbf{答案：}$0.9$

\textbf{分析：}设$OD_1=DC_1=CB_1=BA_1=1$，则可得关于$k_3$的方程，求出其解后可得正确的选项.

\textbf{详解：}设$OD_1=DC_1=CB_1=BA_1=1$，则$CC_1=k_1,BB_1=k_2,AA_1=k_3$，依题意，有$k_3-0.2=k_1,k_3-0.1=k_2$，且$\dfrac{DD_1+CC_1+BB_1+AA_1}{OD_1+DC_1+CB_1+BA_1}=0.725$，所以$\dfrac{0.5+3k_3-0.3}{4}=0.725$，故$k_3=0.9$，故选D.
\end{solution}
\end{question}

\begin{question}[points=5]
已知向量 $\vec{a}=(3,4),\vec{b}=(1,0),\vec{c}=\vec{a}+t\vec{b}$，若 $\langle\vec{a},\vec{c}\rangle=\langle\vec{b},\vec{c}\rangle$，则 $t=$（ ）
\begin{choices}
  \item $-6$
  \item $-5$
  \item $5$
  \item $6$
\end{choices}
\begin{solution}
\textbf{答案：}$5$

\textbf{分析：}利用向量的运算和向量的夹角的余弦公式的坐标形式化简即可求得.

\textbf{详解：}$\vec{c}=(3+t,4)$，$\cos\langle\vec{a},\vec{c}\rangle=\cos\langle\vec{b},\vec{c}\rangle$，即$\dfrac{9+3t+16}{5|c|}=\dfrac{3+t}{|c|}$，解得$t=5$，故选C.
\end{solution}
\end{question}

\begin{question}[points=5]
有甲、乙、丙、丁、戊5名同学站成一排参加文艺汇演，若甲不站在两端，丙和丁相邻，则不同排列方式共有（ ）
\begin{choices}
  \item 12种
  \item 24种
  \item 36种
  \item 48种
\end{choices}
\begin{solution}
\textbf{答案：}24种

\textbf{分析：}利用捆绑法处理丙丁，用插空法安排甲，利用排列组合与计数原理即可得解.

\textbf{详解：}因为丙丁要在一起，先把丙丁捆绑，看做一个元素，连同乙，戊看成三个元素排列，有$3!$种排列方式；为使甲不在两端，必须且只需甲在此三个元素的中间两个位置任选一个位置插入，有2种插空方式；注意到丙丁两人的顺序可交换，有2种排列方式，故安排这5名同学共有$3!\times 2\times 2=24$种不同的排列方式，故选B.
\end{solution}
\end{question}

\begin{question}[points=5]
若 $\sin(\alpha+\beta)+\cos(\alpha+\beta)=2\sqrt{2}\cos\left(\alpha+\dfrac{\pi}{4}\right)\sin\beta$，则（ ）
\begin{choices}
  \item $\tan(\alpha-\beta)=1$
  \item $\tan(\alpha+\beta)=1$
  \item $\tan(\alpha-\beta)=-1$
  \item $\tan(\alpha+\beta)=-1$
\end{choices}
\begin{solution}
\textbf{答案：}$\tan(\alpha-\beta)=-1$

\textbf{分析：}由两角和差的正余弦公式化简，结合同角三角函数的商数关系即可得解.

\textbf{详解：}由已知得：$\sin\alpha\cos\beta+\cos\alpha\sin\beta+\cos\alpha\cos\beta-\sin\alpha\sin\beta=2(\cos\alpha-\sin\alpha)\sin\beta$，即$\sin\alpha\cos\beta-\cos\alpha\sin\beta+\cos\alpha\cos\beta+\sin\alpha\sin\beta=0$，即$\sin(\alpha-\beta)+\cos(\alpha-\beta)=0$，所以$\tan(\alpha-\beta)=-1$，故选C.
\end{solution}
\end{question}

\begin{question}[points=5]
已知正三棱台的高为1，上、下底面边长分别为$3\sqrt{3}$和$4\sqrt{3}$，其顶点都在同一球面上，则该球的表面积为（ ）
\begin{choices}
  \item $100\pi$
  \item $128\pi$
  \item $144\pi$
  \item $192\pi$
\end{choices}
\begin{solution}
\textbf{答案：}$100\pi$

\textbf{分析：}根据题意可求出正三棱台上下底面所在圆面的半径$r_1,r_2$，再根据球心距，圆面半径，以及球的半径之间的关系，即可解出球的半径，从而得出球的表面积.

\textbf{详解：}设正三棱台上下底面所在圆面的半径$r_1,r_2$，所以$2r_1=\dfrac{3\sqrt{3}}{\sin60^\circ},2r_2=\dfrac{4\sqrt{3}}{\sin60^\circ}$，即$r_1=3,r_2=4$，设球心到上下底面的距离分别为$d_1,d_2$，球的半径为$R$，所以$d_1=\sqrt{R^2-9},d_2=\sqrt{R^2-16}$，故$|d_1-d_2|=1$，即$\sqrt{R^2-9}-\sqrt{R^2-16}=1$或$\sqrt{R^2-9}+\sqrt{R^2-16}=1$，解得$R^2=25$符合题意，所以球的表面积为$S=4\pi R^2=100\pi$，故选A.
\end{solution}
\end{question}

\begin{question}[points=5]
已知函数$f(x)$的定义域为$\mathbf{R}$，且$f(x+y)+f(x-y)=f(x)f(y),f(1)=1$，则$\displaystyle\sum_{k=1}^{22}f(k)=$（ ）
\begin{choices}
  \item $-3$
  \item $-2$
  \item $0$
  \item $1$
\end{choices}
\begin{solution}
\textbf{答案：}$-3$

\textbf{分析：}根据题意赋值即可知函数$f(x)$的一个周期为6，求出函数一个周期中的$f(1),f(2),\cdots,f(6)$的值，即可解出.

\textbf{详解：}因为$f(x+y)+f(x-y)=f(x)f(y)$，令$x=1,y=0$可得$2f(1)=f(1)f(0)$，所以$f(0)=2$，令$x=0$可得$f(y)+f(-y)=2f(y)$，即$f(y)=f(-y)$，所以函数$f(x)$为偶函数，令$y=1$得$f(x+1)+f(x-1)=f(x)f(1)=f(x)$，即有$f(x+2)+f(x)=f(x+1)$，从而可知$f(x+2)=-f(x-1)$，$f(x-1)=-f(x-4)$，故$f(x+2)=f(x-4)$，即$f(x)=f(x+6)$，所以函数$f(x)$的一个周期为6. 因为$f(2)=f(1)-f(0)=1-2=-1$，$f(3)=f(2)-f(1)=-1-1=-2$，$f(4)=f(-2)=f(2)=-1$，$f(5)=f(-1)=f(1)=1$，$f(6)=f(0)=2$，所以一个周期内的$f(1)+f(2)+\cdots+f(6)=0$．由于22除以6余4，所以$\sum_{k=1}^{22}f(k)=f(1)+f(2)+f(3)+f(4)=1-1-2-1=-3$，故选A.
\end{solution}
\end{question}

\section*{二、多项选择题}
本题共4小题，每小题5分，共20分．在每小题给出的选项中，有多项符合题目要求．全部选对的得5分，部分选对的得2分，有选错的得0分.

\begin{question}[points=5]
已知函数$f(x)=\sin(2x+\varphi)(0<\varphi<\pi)$的图像关于点$\left(\dfrac{2\pi}{3},0\right)$中心对称，则（ ）
\begin{choices}
  \item $f(x)$在区间$\left(0,\dfrac{5\pi}{12}\right)$单调递减
  \item $f(x)$在区间$\left(-\dfrac{\pi}{12},\dfrac{11\pi}{12}\right)$有两个极值点
  \item 直线$x=\dfrac{7\pi}{6}$是曲线$y=f(x)$的对称轴
  \item 直线$y=\dfrac{\sqrt{3}}{2}-x$是曲线$y=f(x)$的切线
\end{choices}
\begin{solution}
\textbf{答案：}AD

\textbf{分析：}根据三角函数的性质逐个判断各选项，即可解出.

\textbf{详解：}由题意得$f\left(\dfrac{2\pi}{3}\right)=\sin\left(\dfrac{4\pi}{3}+\varphi\right)=0$，所以$\dfrac{4\pi}{3}+\varphi=k\pi,k\in\mathbf{Z}$，即$\varphi=-\dfrac{4\pi}{3}+k\pi,k\in\mathbf{Z}$，又$0<\varphi<\pi$，所以$k=2$时，$\varphi=\dfrac{2\pi}{3}$，故$f(x)=\sin\left(2x+\dfrac{2\pi}{3}\right)$．对A，当$x\in\left(0,\dfrac{5\pi}{12}\right)$时，$2x+\dfrac{2\pi}{3}\in\left(\dfrac{2\pi}{3},\dfrac{3\pi}{2}\right)$，由正弦函数$y=\sin u$图象知$y=f(x)$在$\left(0,\dfrac{5\pi}{12}\right)$上是单调递减；对B，当$x\in\left(-\dfrac{\pi}{12},\dfrac{11\pi}{12}\right)$时，$2x+\dfrac{2\pi}{3}\in\left(\dfrac{\pi}{2},\dfrac{5\pi}{2}\right)$，由正弦函数$y=\sin u$图象知$y=f(x)$只有1个极值点，由$2x+\dfrac{2\pi}{3}=\dfrac{3\pi}{2}$，解得$x=\dfrac{5\pi}{12}$，即$x=\dfrac{5\pi}{12}$为函数的唯一极值点；对C，当$x=\dfrac{7\pi}{6}$时，$2x+\dfrac{2\pi}{3}=3\pi$，$f\left(\dfrac{7\pi}{6}\right)=0$，直线$x=\dfrac{7\pi}{6}$不是对称轴；对D，由$y'=2\cos\left(2x+\dfrac{2\pi}{3}\right)=-1$得$\cos\left(2x+\dfrac{2\pi}{3}\right)=-\dfrac{1}{2}$，解得$2x+\dfrac{2\pi}{3}=\dfrac{2\pi}{3}+2k\pi$或$2x+\dfrac{2\pi}{3}=\dfrac{4\pi}{3}+2k\pi,k\in\mathbf{Z}$，从而得$x=k\pi$或$x=\dfrac{\pi}{3}+k\pi,k\in\mathbf{Z}$，所以函数$y=f(x)$在点$\left(0,\dfrac{\sqrt{3}}{2}\right)$处的切线斜率为$k=y'|_{x=0}=2\cos\dfrac{2\pi}{3}=-1$，切线方程为$y-\dfrac{\sqrt{3}}{2}=-(x-0)$即$y=\dfrac{\sqrt{3}}{2}-x$．故选AD.
\end{solution}
\end{question}

\begin{question}[points=5]
已知$O$为坐标原点，过抛物线$C:y^2=2px(p>0)$焦点$F$的直线与$C$交于$A,B$两点，其中$A$在第一象限，点$M(p,0)$，若$|AF|=|AM|$，则（ ）
\begin{choices}
  \item 直线$AB$的斜率为$2\sqrt{6}$
  \item $|OB|=|OF|$
  \item $|AB|>4|OF|$
  \item $\angle OAM+\angle OBM<180^\circ$
\end{choices}
\begin{solution}
\textbf{答案：}ACD

\textbf{分析：}由$AF=AM$及抛物线方程求得$A\left(\dfrac{3p}{4},\dfrac{\sqrt{6}p}{2}\right)$，再由斜率公式即可判断A选项；表示出直线$AB$的方程，联立抛物线求得$B\left(\dfrac{p}{3},-\dfrac{\sqrt{6}p}{3}\right)$，即可求出$OB$判断B选项；由抛物线的定义求出$AB=\dfrac{25p}{12}$即可判断C选项；由$\overrightarrow{OA}\cdot\overrightarrow{OB}<0$，$\overrightarrow{MA}\cdot\overrightarrow{MB}<0$求得$\angle AOB$，$\angle AMB$为钝角即可判断D选项.

\textbf{详解：}对于A，易得$F\left(\dfrac{p}{2},0\right)$，由$AF=AM$可得点$A$在$FM$的垂直平分线上，则$A$横坐标为$\dfrac{\frac{p}{2}+p}{2}=\dfrac{3p}{4}$，代入抛物线可得$y^2=2p\cdot\dfrac{3p}{4}=\dfrac{3p^2}{2}$，则$A\left(\dfrac{3p}{4},\dfrac{\sqrt{6}p}{2}\right)$，则直线$AB$的斜率为$\dfrac{\frac{\sqrt{6}p}{2}}{\frac{3p}{4}-\frac{p}{2}}=2\sqrt{6}$，A正确；对于B，由斜率为$2\sqrt{6}$可得直线$AB$的方程为$x=\dfrac{1}{2\sqrt{6}}y+\dfrac{p}{2}$，联立抛物线方程得$y^2-\dfrac{\sqrt{6}}{6}py-p^2=0$，设$B(x_1,y_1)$，则$\dfrac{\sqrt{6}}{2}p+y_1=\dfrac{\sqrt{6}}{6}p$，则$y_1=-\dfrac{\sqrt{6}}{3}p$，代入抛物线得$\left(-\dfrac{\sqrt{6}}{3}p\right)^2=2px_1$，解得$x_1=\dfrac{p}{3}$，则$B\left(\dfrac{p}{3},-\dfrac{\sqrt{6}p}{3}\right)$，则$|OB|=\sqrt{\left(\dfrac{p}{3}\right)^2+\left(-\dfrac{\sqrt{6}p}{3}\right)^2}=\dfrac{\sqrt{7}p}{3}\neq|OF|=\dfrac{p}{2}$，B错误；对于C，由抛物线定义知：$|AB|=\dfrac{3p}{4}+\dfrac{p}{3}+p=\dfrac{25p}{12}>2p=4|OF|$，C正确；对于D，$\overrightarrow{OA}\cdot\overrightarrow{OB}=\left(\dfrac{3p}{4},\dfrac{\sqrt{6}p}{2}\right)\cdot\left(\dfrac{p}{3},-\dfrac{\sqrt{6}p}{3}\right)=\dfrac{3p}{4}\cdot\dfrac{p}{3}+\dfrac{\sqrt{6}p}{2}\cdot\left(-\dfrac{\sqrt{6}p}{3}\right)=\dfrac{p^2}{4}-\dfrac{p^2}{2}=-\dfrac{p^2}{4}<0$，则$\angle AOB$为钝角，又$\overrightarrow{MA}\cdot\overrightarrow{MB}=\left(-\dfrac{p}{4},\dfrac{\sqrt{6}p}{2}\right)\cdot\left(-\dfrac{2p}{3},-\dfrac{\sqrt{6}p}{3}\right)=\left(-\dfrac{p}{4}\right)\cdot\left(-\dfrac{2p}{3}\right)+\dfrac{\sqrt{6}p}{2}\cdot\left(-\dfrac{\sqrt{6}p}{3}\right)=\dfrac{p^2}{6}-p^2=-\dfrac{5p^2}{6}<0$，则$\angle AMB$为钝角，又$\angle AOB+\angle AMB+\angle OAM+\angle OBM=360^\circ$，则$\angle OAM+\angle OBM<180^\circ$，D正确. 故选ACD.
\end{solution}
\end{question}

\begin{question}[points=5]
如图，四边形$ABCD$为正方形，$ED\perp$平面$ABCD$，$FB\parallel ED$，$AB=ED=2FB$，记三棱锥$E-ACD$，$F-ABC$，$F-ACE$的体积分别为$V_1,V_2,V_3$，则（ ）

\begin{center}
  \includegraphics[width=.38\linewidth]{question-11-1.jpg}
\end{center}
\begin{choices}
  \item $V_3=2V_2$
  \item $V_3=V_1$
  \item $V_3=V_1+V_2$
  \item $2V_3=3V_1$
\end{choices}
\begin{solution}
\textbf{答案：}CD

\textbf{分析：}直接由体积公式计算$V_1,V_2$，连接$BD$交$AC$于点$M$，连接$EM,FM$，由$V_3=V_{A-EFM}+V_{C-EFM}$计算出$V_3$，依次判断选项即可.

\textbf{详解：}设$AB=ED=2FB=2a$，因为$ED\perp$平面$ABCD$，$FB\parallel ED$，则$V_1=\dfrac{1}{3}\cdot ED\cdot S_{\triangle ACD}=\dfrac{1}{3}\cdot 2a\cdot \dfrac{1}{2}\cdot(2a)^2=\dfrac{4}{3}a^3$，$V_2=\dfrac{1}{3}\cdot FB\cdot S_{\triangle ABC}=\dfrac{1}{3}\cdot a\cdot \dfrac{1}{2}\cdot (2a)^2=\dfrac{2}{3}a^3$，连接$BD$交$AC$于点$M$，连接$EM,FM$，易得$BD\perp AC$，又$ED\perp$平面$ABCD$，$AC\subset$平面$ABCD$，则$ED\perp AC$，又$ED\cap BD=D$，$ED,BD\subset$平面$BDEF$，则$AC\perp$平面$BDEF$，又$BM=DM=\dfrac{1}{2}BD=\sqrt{2}a$，过$F$作$FG\perp DE$于$G$，易得四边形$BDGF$为矩形，则$FG=BD=2\sqrt{2}a,EG=a$，则$EM=\sqrt{(\sqrt{2}a)^2+(\sqrt{2}a)^2}=2a,FM=\sqrt{a^2+(\sqrt{2}a)^2}=\sqrt{3}a,EF=\sqrt{a^2+(2\sqrt{2}a)^2}=3a$，$EM^2+FM^2=EF^2$，则$EM\perp FM$，$S_{\triangle EFM}=\dfrac{1}{2}EM\cdot FM=\dfrac{1}{2}\cdot 2a\cdot\sqrt{3}a=\sqrt{3}a^2$，$AC=2\sqrt{2}a$，则$V_3=V_{A-EFM}+V_{C-EFM}=\dfrac{1}{3}AC\cdot S_{\triangle EFM}=2a^3$，则$2V_3=3V_1$，$V_3=3V_2$，$V_3=V_1+V_2$，故AB错误；CD正确. 故选CD.
\end{solution}
\end{question}

\begin{question}[points=5]
若$x,y$满足$x^2+y^2-xy=1$，则（ ）
\begin{choices}
  \item $x+y\leq 1$
  \item $x+y\geq -2$
  \item $x^2+y^2\leq 2$
  \item $x^2+y^2\geq 1$
\end{choices}
\begin{solution}
\textbf{答案：}BC

\textbf{分析：}根据基本不等式或者取特值即可判断各选项的真假.

\textbf{详解：}因为$ab\leq\left(\dfrac{a+b}{2}\right)^2\leq\dfrac{a^2+b^2}{2}$（$a,b\in\mathbf{R}$），由$x^2+y^2-xy=1$可变形为$(x+y)^2-1=3xy\leq 3\left(\dfrac{x+y}{2}\right)^2$，解得$-2\leq x+y\leq 2$，当且仅当$x=y=-1$时，$x+y=-2$，当且仅当$x=y=1$时，$x+y=2$，所以A错误，B正确；由$x^2+y^2-xy=1$可变形为$x^2+y^2-1=xy\leq\dfrac{x^2+y^2}{2}$，解得$x^2+y^2\leq 2$，当且仅当$x=y=\pm1$时取等号，所以C正确；因为$x^2+y^2-xy=1$变形可得$\left(x-\dfrac{y}{2}\right)^2+\dfrac{3}{4}y^2=1$，设$x-\dfrac{y}{2}=\cos\theta,\dfrac{\sqrt{3}}{2}y=\sin\theta$，所以$x=\cos\theta+\dfrac{1}{\sqrt{3}}\sin\theta,y=\dfrac{2}{\sqrt{3}}\sin\theta$，因此$x^2+y^2=\cos^2\theta+\dfrac{5}{3}\sin^2\theta+\dfrac{2}{\sqrt{3}}\sin\theta\cos\theta=1+\dfrac{1}{\sqrt{3}}\sin2\theta-\dfrac{1}{3}\cos2\theta+\dfrac{1}{3}=\dfrac{4}{3}+\dfrac{2}{3}\sin\left(2\theta-\dfrac{\pi}{6}\right)\in\left[\dfrac{2}{3},2\right]$，所以当$x=\dfrac{\sqrt{3}}{3},y=-\dfrac{\sqrt{3}}{3}$时满足等式，但是$x^2+y^2\geq 1$不成立，所以D错误．故选BC.
\end{solution}
\end{question}

\section*{三、填空题}
本题共4小题，每小题5分，共20分.

\begin{question}[points=5]
已知随机变量$X$服从正态分布$N(2,\sigma^2)$，且$P(2<X\leq 2.5)=0.36$，则$P(X>2.5)=$.
\fillin[{$0.14$}]
\begin{solution}
\textbf{答案：}$0.14$

\textbf{分析：}根据正态分布曲线的性质即可解出.

\textbf{详解：}因为$X\sim N(2,\sigma^2)$，所以$P(X<2)=P(X>2)=0.5$，因此$P(X>2.5)=P(X>2)-P(2<X\leq 2.5)=0.5-0.36=0.14$．故答案为$0.14$．
\end{solution}
\end{question}

\begin{question}[points=5]
曲线$y=\ln|x|$过坐标原点的两条切线的方程为，.
\fillin[{$y=\dfrac{1}{e}x$，$y=-\dfrac{1}{e}x$}]
\begin{solution}
\textbf{答案：}$y=\dfrac{1}{e}x$，$y=-\dfrac{1}{e}x$

\textbf{分析：}分$x>0$和$x<0$两种情况，当$x>0$时设切点为$(x_0,\ln x_0)$，求出函数的导函数，即可求出切线的斜率，从而表示出切线方程，再根据切线过坐标原点求出$x_0$，即可求出切线方程，当$x<0$时同理可得.

\textbf{详解：}解：因为$y=\ln|x|$，当$x>0$时$y=\ln x$，设切点为$(x_0,\ln x_0)$，由$y'=\dfrac{1}{x}$，所以$y'|_{x=x_0}=\dfrac{1}{x_0}$，所以切线方程为$y-\ln x_0=\dfrac{1}{x_0}(x-x_0)$，又切线过坐标原点，所以$-\ln x_0=\dfrac{1}{x_0}(-x_0)$，解得$x_0=e$，所以切线方程为$y-1=\dfrac{1}{e}(x-e)$，即$y=\dfrac{1}{e}x$；当$x<0$时$y=\ln(-x)$，设切点为$(x_1,\ln(-x_1))$，由$y'=\dfrac{1}{x}$（因为$(-x)'=-1$，复合函数求导），得$y'=\dfrac{1}{x}$（注意此处$x<0$，$\dfrac{1}{x}$也为负），所以$y'|_{x=x_1}=\dfrac{1}{x_1}$，所以切线方程为$y-\ln(-x_1)=\dfrac{1}{x_1}(x-x_1)$，又切线过坐标原点，所以$-\ln(-x_1)=\dfrac{1}{x_1}(-x_1)$，解得$x_1=-e$，所以切线方程为$y-1=\dfrac{1}{-e}(x+e)$，即$y=-\dfrac{1}{e}x$；故答案为$y=\dfrac{1}{e}x$；$y=-\dfrac{1}{e}x$．
\end{solution}
\end{question}

\begin{question}[points=5]
设点$A(-2,3),B(0,a)$，若直线$AB$关于$y=a$对称的直线与圆$(x+3)^2+(y+2)^2=1$有公共点，则$a$的取值范围是.
\fillin[{$\left[\dfrac{1}{3},\dfrac{3}{2}\right]$}]
\begin{solution}
\textbf{答案：}$\left[\dfrac{1}{3},\dfrac{3}{2}\right]$

\textbf{分析：}首先求出点$A$关于$y=a$对称点$A'$的坐标，即可得到直线$l$的方程，根据圆心到直线的距离小于等于半径得到不等式，解得即可.

\textbf{详解：}解：$A(-2,3)$关于$y=a$对称的点的坐标为$A'(-2,2a-3)$，$B(0,a)$在直线$y=a$上，所以$A'B$所在直线即为直线$l$，所以直线$l$为$y=\dfrac{a-3}{-2}x+a$，即$(a-3)x+2y-2a=0$；圆$C:(x+3)^2+(y+2)^2=1$，圆心$C(-3,-2)$，半径$r=1$，依题意圆心到直线$l$的距离$d=\dfrac{|-3(a-3)-4-2a|}{\sqrt{(a-3)^2+2^2}}\leq 1$，即$|5-5a|\leq\sqrt{(a-3)^2+4}$，两边平方得$(5-5a)^2\leq (a-3)^2+4$，整理得$24a^2-44a+12\leq 0$，即$6a^2-11a+3\leq 0$，解得$\dfrac{1}{3}\leq a\leq \dfrac{3}{2}$，所以$a$的取值范围是$\left[\dfrac{1}{3},\dfrac{3}{2}\right]$．故答案为$\left[\dfrac{1}{3},\dfrac{3}{2}\right]$.
\end{solution}
\end{question}

\begin{question}[points=5]
已知直线$l$与椭圆$\dfrac{x^2}{6}+\dfrac{y^2}{3}=1$在第一象限交于$A,B$两点，$l$与$x$轴，$y$轴分别交于$M,N$两点，且$|MA|=|NB|,|MN|=2\sqrt{3}$，则$l$的方程为.
\fillin[{$x+\sqrt{2}y-2\sqrt{2}=0$}]
\begin{solution}
\textbf{答案：}$x+\sqrt{2}y-2\sqrt{2}=0$

\textbf{分析：}令$AB$的中点为$E$，设$A(x_1,y_1),B(x_2,y_2)$，利用点差法得到$k_{OE}\cdot k_{AB}=-\dfrac{1}{2}$，设直线$AB:y=kx+m,k<0,m>0$，求出$M,N$的坐标，再根据$MN$求出$k,m$，即可得解.

\textbf{详解：}解：令$AB$的中点为$E$，因为$|MA|=|NB|$，所以$|ME|=|NE|$，设$A(x_1,y_1),B(x_2,y_2)$，则$\dfrac{x_1^2}{6}+\dfrac{y_1^2}{3}=1$，$\dfrac{x_2^2}{6}+\dfrac{y_2^2}{3}=1$，两式相减得$\dfrac{(x_1-x_2)(x_1+x_2)}{6}+\dfrac{(y_1-y_2)(y_1+y_2)}{3}=0$，所以$\dfrac{y_1+y_2}{x_1+x_2}\cdot\dfrac{y_1-y_2}{x_1-x_2}=-\dfrac{1}{2}$，即$k_{OE}\cdot k_{AB}=-\dfrac{1}{2}$，设直线$AB:y=kx+m,k<0,m>0$，令$x=0$得$y=m$，令$y=0$得$x=-\dfrac{m}{k}$，即$M\left(-\dfrac{m}{k},0\right),N(0,m)$，所以$E\left(-\dfrac{m}{2k},\dfrac{m}{2}\right)$，则$k_{OE}=\dfrac{m/2}{-m/(2k)}=-k$，所以$(-k)\cdot k=-\dfrac{1}{2}$，解得$k^2=\dfrac{1}{2}$，由于$k<0$，所以$k=-\dfrac{\sqrt{2}}{2}$，又$|MN|=\sqrt{m^2+\left(\dfrac{m}{k}\right)^2}=2\sqrt{3}$，即$\sqrt{m^2+(\sqrt{2}m)^2}=2\sqrt{3}$，解得$m=2$（$m>0$），所以直线$AB$的方程为$y=-\dfrac{\sqrt{2}}{2}x+2$，即$x+\sqrt{2}y-2\sqrt{2}=0$．故答案为$x+\sqrt{2}y-2\sqrt{2}=0$.
\end{solution}
\end{question}

\section*{四、解答题}
本题共6小题，共70分．解答应写出文字说明，证明过程或演算步骤.

\begin{question}[points=10]
已知$\{a_n\}$为等差数列，$\{b_n\}$是公比为2的等比数列，且$a_2-b_2=a_3-b_3=b_4-a_4$．

(1)证明：$a_1=b_1$；

(2)求集合$\{k|b_k=a_m+a_1,1\leq m\leq 500\}$中元素个数．
\begin{solution}
\textbf{答案：}(1)证明见解析；(2)9

\textbf{分析：}(1)设数列$\{a_n\}$的公差为$d$，根据题意列出方程组即可证出；(2)根据题意化简可得$m=2^{k-2}$，即可解出.

\textbf{详解：}（1）设数列$\{a_n\}$的公差为$d$，所以$\begin{cases} a_1+d-2b_1=a_1+2d-4b_1 \\ a_1+d-2b_1=8b_1-(a_1+3d) \end{cases}$，即可解得$b_1=a_1=\dfrac{d}{2}$，所以原命题得证．
（2）由（1）知，$b_1=a_1=\dfrac{d}{2}$，所以$b_k=a_m+a_1\Leftrightarrow b_1\times 2^{k-1}=a_1+(m-1)d+a_1$，即$2^{k-1}=2m$，亦即$m=2^{k-2}\in[1,500]$，解得$2\leq k\leq 10$，所以满足等式的解$k=2,3,4,\cdots,10$，故集合$\{k|b_k=a_m+a_1,1\leq m\leq 500\}$中的元素个数为$10-2+1=9$.
\end{solution}
\end{question}

\begin{question}[points=12]
记$\triangle ABC$的内角$A,B,C$的对边分别为$a,b,c$，分别以$a,b,c$为边长的三个正三角形的面积依次为$S_1,S_2,S_3$，已知$S_1-S_2+S_3=\dfrac{\sqrt{3}}{2},\sin B=\dfrac{1}{3}$．

(1)求$\triangle ABC$的面积；

(2)若$\sin A\sin C=\dfrac{\sqrt{2}}{3}$，求$b$．
\begin{solution}
\textbf{答案：}(1)$\dfrac{\sqrt{2}}{8}$；(2)$\dfrac{1}{2}$

\textbf{分析：}(1)先表示出$S_1,S_2,S_3$，再由$S_1-S_2+S_3=\dfrac{\sqrt{3}}{2}$求得$a^2+c^2-b^2=2$，结合余弦定理及平方关系求得$ac$，再由面积公式求解即可；(2)由正弦定理得$\dfrac{b^2}{\sin^2 B}=\dfrac{ac}{\sin A\sin C}$，即可求解.

\textbf{详解：}（1）由题意得$S_1=\dfrac{1}{2}\cdot a^2\cdot\dfrac{\sqrt{3}}{2}=\dfrac{\sqrt{3}}{4}a^2$，$S_2=\dfrac{\sqrt{3}}{4}b^2$，$S_3=\dfrac{\sqrt{3}}{4}c^2$，则$S_1-S_2+S_3=\dfrac{\sqrt{3}}{4}a^2-\dfrac{\sqrt{3}}{4}b^2+\dfrac{\sqrt{3}}{4}c^2=\dfrac{\sqrt{3}}{2}$，即$a^2+c^2-b^2=2$，由余弦定理得$\cos B=\dfrac{a^2+c^2-b^2}{2ac}$，整理得$ac\cos B=1$，则$\cos B>0$，又$\sin B=\dfrac{1}{3}$，则$\cos B=\sqrt{1-\left(\dfrac{1}{3}\right)^2}=\dfrac{2\sqrt{2}}{3}$，$ac=\dfrac{1}{\cos B}=\dfrac{3}{2\sqrt{2}}$，则$S_{\triangle ABC}=\dfrac{1}{2}ac\sin B=\dfrac{\sqrt{2}}{8}$．
（2）由正弦定理得$\dfrac{b}{\sin B}=\dfrac{a}{\sin A}=\dfrac{c}{\sin C}$，则$\dfrac{b^2}{\sin^2 B}=\dfrac{a}{\sin A}\cdot\dfrac{c}{\sin C}=\dfrac{ac}{\sin A\sin C}=\dfrac{\frac{3}{2\sqrt{2}}}{\frac{\sqrt{2}}{3}}=\dfrac{9}{4}$，则$\dfrac{b}{\sin B}=\dfrac{3}{2}$，$b=\dfrac{3}{2}\sin B=\dfrac{1}{2}$.
\end{solution}
\end{question}

\begin{question}[points=12]
在某地区进行流行病学调查，随机调查了100位某种疾病患者的年龄，得到如下的样本数据的频率分布直方图：

(1)估计该地区这种疾病患者的平均年龄（同一组中的数据用该组区间的中点值为代表）；

(2)估计该地区一位这种疾病患者的年龄位于区间$[20,70)$的概率；

(3)已知该地区这种疾病的患病率为$0.1\%$，该地区年龄位于区间$[40,50)$的人口占该地区总人口的$16\%$．从该地区中任选一人，若此人的年龄位于区间$[40,50)$，求此人患这种疾病的概率．（以样本数据中患者的年龄位于各区间的频率作为患者的年龄位于该区间的概率，精确到$0.0001$）.

\begin{center}
  \includegraphics[width=.38\linewidth]{question-19-1.jpg}
\end{center}
\begin{solution}
\textbf{答案：}(1)$47.9$岁；(2)$0.89$；(3)$0.0014$

\textbf{分析：}(1)根据平均值等于各矩形的面积乘以对应区间的中点值的和即可求出；(2)设$A=\{$一人患这种疾病的年龄在区间$[20,70)\}$，根据对立事件的概率公式$P(A)=1-P(\overline{A})$即可解出；(3)根据条件概率公式即可求出.

\textbf{详解：}（1）平均年龄$\bar{x}=(5\times0.001+15\times0.002+25\times0.012+35\times0.017+45\times0.023+55\times0.020+65\times0.017+75\times0.006+85\times0.002)\times10=47.9$（岁）．
（2）设$A=\{$一人患这种疾病的年龄在区间$[20,70)\}$，所以$P(A)=1-P(\overline{A})=1-(0.001+0.002+0.006+0.002)\times10=1-0.11=0.89$．
（3）设$B=$“任选一人年龄位于区间$[40,50)$”，$C=$“从该地区中任选一人患这种疾病”，则由已知得：$P(B)=16\%=0.16,P(C)=0.1\%=0.001,P(B|C)=0.023\times10=0.23$，则由条件概率公式可得从该地区中任选一人，若此人的年龄位于区间$[40,50)$，此人患这种疾病的概率为$P(C|B)=\dfrac{P(BC)}{P(B)}=\dfrac{P(C)P(B|C)}{P(B)}=\dfrac{0.001\times0.23}{0.16}=0.0014375\approx 0.0014$．
\end{solution}
\end{question}

\begin{question}[points=12]
如图，$PO$是三棱锥$P-ABC$的高，$PA=PB$，$AB\perp AC$，$E$是$PB$的中点．

(1)证明：$OE\parallel$平面$PAC$；

(2)若$\angle ABO=\angle CBO=30^\circ$，$PO=3$，$PA=5$，求二面角$C-AE-B$的正弦值．

\begin{center}
  \includegraphics[width=.38\linewidth]{question-20-1.jpg}
\end{center}
\begin{solution}
\textbf{答案：}（1）证明见解析；（2）$\dfrac{11}{13}$

\textbf{分析：}（1）连接$BO$并延长交$AC$于点$D$，连接$OA,PD$，根据三角形全等得到$OA=OB$，再根据直角三角形的性质得到$AO=DO$，即可得到$O$为$BD$的中点从而得到$OE\parallel PD$，即可得证；（2）建立适当的空间直角坐标系，利用空间向量法求出二面角的余弦的绝对值，再根据同角三角函数的基本关系计算可得.

\textbf{详解：}（1）证明：连接$BO$并延长交$AC$于点$D$，连接$OA,PD$，因为$PO$是三棱锥$P-ABC$的高，所以$PO\perp$平面$ABC$，$AO,BO\subset$平面$ABC$，所以$PO\perp AO,PO\perp BO$，又$PA=PB$，所以$\triangle POA\cong\triangle POB$，即$OA=OB$，所以$\angle OAB=\angle OBA$，又$AB\perp AC$，即$\angle BAC=90^\circ$，所以$\angle OAB+\angle OAD=90^\circ$，$\angle OBA+\angle ODA=90^\circ$，所以$\angle ODA=\angle OAD$，所以$AO=DO$，即$AO=DO=OB$，所以$O$为$BD$的中点，又$E$为$PB$的中点，所以$OE\parallel PD$，又$OE\not\subset$平面$PAC$，$PD\subset$平面$PAC$，所以$OE\parallel$平面$PAC$．
（2）解：过点$A$作$Az\parallel OP$，如图建立平面直角坐标系，因为$PO=3$，$AP=5$，所以$OA=\sqrt{AP^2-PO^2}=4$，又$\angle OBA=\angle OBC=30^\circ$，所以$BD=2OA=8$，则$AD=4$，$AB=4\sqrt{3}$，所以$AC=12$，所以$O(2\sqrt{3},2,0)$，$B(4\sqrt{3},0,0)$，$P(2\sqrt{3},2,3)$，$C(0,12,0)$，所以$E(3\sqrt{3},1,\dfrac{3}{2})$，则$\overrightarrow{AE}=(3\sqrt{3},1,\dfrac{3}{2})$，$\overrightarrow{AB}=(4\sqrt{3},0,0)$，$\overrightarrow{AC}=(0,12,0)$，设平面$AEB$的法向量为$\vec{n}=(x,y,z)$，则$\begin{cases} \vec{n}\cdot\overrightarrow{AE}=3\sqrt{3}x+y+\dfrac{3}{2}z=0 \\ \vec{n}\cdot\overrightarrow{AB}=4\sqrt{3}x=0 \end{cases}$，令$z=2$，则$y=-3$，$x=0$，所以$\vec{n}=(0,-3,2)$；设平面$AEC$的法向量为$\vec{m}=(a,b,c)$，则$\begin{cases} \vec{m}\cdot\overrightarrow{AE}=3\sqrt{3}a+b+\dfrac{3}{2}c=0 \\ \vec{m}\cdot\overrightarrow{AC}=12b=0 \end{cases}$，令$a=\sqrt{3}$，则$c=-6$，$b=0$，所以$\vec{m}=(\sqrt{3},0,-6)$；所以$\cos\langle\vec{n},\vec{m}\rangle=\dfrac{\vec{n}\cdot\vec{m}}{|\vec{n}||\vec{m}|}=\dfrac{-12}{\sqrt{13}\times\sqrt{39}}=-\dfrac{4\sqrt{3}}{13}$．设二面角$C-AE-B$的大小为$\theta$，则$\cos\theta=|\cos\langle\vec{n},\vec{m}\rangle|=\dfrac{4\sqrt{3}}{13}$，所以$\sin\theta=\sqrt{1-\cos^2\theta}=\sqrt{1-\left(\dfrac{4\sqrt{3}}{13}\right)^2}=\dfrac{11}{13}$，即二面角$C-AE-B$的正弦值为$\dfrac{11}{13}$.
\end{solution}
\end{question}

\begin{question}[points=12]
已知双曲线$C:\dfrac{x^2}{a^2}-\dfrac{y^2}{b^2}=1(a>0,b>0)$的右焦点为$F(2,0)$，渐近线方程为$y=\pm\sqrt{3}x$．

(1)求$C$的方程；

(2)过$F$的直线与$C$的两条渐近线分别交于$A,B$两点，点$P(x_1,y_1),Q(x_2,y_2)$在$C$上，且$x_1>x_2>0,y_1>0$．过$P$且斜率为$-\sqrt{3}$的直线与过$Q$且斜率为$\sqrt{3}$的直线交于点$M$．从下面①②③中选取两个作为条件，证明另外一个成立：

①$M$在$AB$上；②$PQ\parallel AB$；③$|MA|=|MB|$．

注：若选择不同的组合分别解答，则按第一个解答计分.
\begin{solution}
\textbf{答案：}（1）$x^2-\dfrac{y^2}{3}=1$；（2）见解析

\textbf{分析：}（1）利用焦点坐标求得$c$的值，利用渐近线方程求得$a,b$的关系，进而利用$a,b,c$的平方关系求得$a,b$的值，得到双曲线的方程；（2）先分析得到直线$AB$的斜率存在且不为零，设直线$AB$的斜率为$k$，$M(x_0,y_0)$，由③$|AM|=|BM|$等价分析得到$x_0+ky_0=\dfrac{8k^2}{k^2-3}$；由直线$PM$和$QM$的斜率得到直线方程，结合双曲线的方程，两点间距离公式得到直线$PQ$的斜率$m=\dfrac{3x_0}{y_0}$，由②$PQ\parallel AB$等价转化为$ky_0=3x_0$，由①$M$在直线$AB$上等价于$ky_0=k^2(x_0-2)$，然后选择两个作为已知条件一个作为结论，进行证明即可.

\textbf{详解：}（1）右焦点为$F(2,0)$，∴$c=2$，∵渐近线方程为$y=\pm\sqrt{3}x$，∴$\dfrac{b}{a}=\sqrt{3}$，∴$b=\sqrt{3}a$，∴$c^2=a^2+b^2=4a^2=4$，∴$a=1$，∴$b=\sqrt{3}$．∴$C$的方程为：$x^2-\dfrac{y^2}{3}=1$．
（2）由已知得直线$PQ$的斜率存在且不为零，直线$AB$的斜率不为零，若选由①②推③或选由②③推①：由②成立可知直线$AB$的斜率存在且不为零；若选①③推②，则$M$为线段$AB$的中点，假若直线$AB$的斜率不存在，则由双曲线的对称性可知$M$在$x$轴上，即为焦点$F$，此时由对称性可知$P$、$Q$关于$x$轴对称，与$x_1=x_2$已知不符；总之，直线$AB$的斜率存在且不为零. 设直线$AB$的斜率为$k$，直线$AB$方程为$y=k(x-2)$，则条件①$M$在$AB$上，等价于$y_0=k(x_0-2)\Leftrightarrow ky_0=k^2(x_0-2)$；两渐近线的方程合并为$3x^2-y^2=0$，联立消去$y$并化简整理得：$(k^2-3)x^2-4k^2x+4k^2=0$，设$A(x_3,y_3),B(x_4,y_4)$，线段中点为$N(x_N,y_N)$，则$x_N=\dfrac{x_3+x_4}{2}=\dfrac{2k^2}{k^2-3}$，$y_N=k(x_N-2)=\dfrac{6k}{k^2-3}$，设$M(x_0,y_0)$，则条件③$|AM|=|BM|$等价于$(x_0-x_3)^2+(y_0-y_3)^2=(x_0-x_4)^2+(y_0-y_4)^2$，移项并利用平方差公式整理得：$(x_3-x_4)[2x_0-(x_3+x_4)]+(y_3-y_4)[2y_0-(y_3+y_4)]=0$，$[2x_0-(x_3+x_4)]+\dfrac{y_3-y_4}{x_3-x_4}[2y_0-(y_3+y_4)]=0$，即$x_0-x_N+k(y_0-y_N)=0$，即$x_0+ky_0=\dfrac{8k^2}{k^2-3}$；由题意知直线$PM$的斜率为$-\sqrt{3}$，直线$QM$的斜率为$\sqrt{3}$，∴由$y_1-y_0=-\sqrt{3}(x_1-x_0)$，$y_2-y_0=\sqrt{3}(x_2-x_0)$，∴$y_1-y_2=-\sqrt{3}(x_1+x_2-2x_0)$，所以直线$PQ$的斜率$m=\dfrac{y_1-y_2}{x_1-x_2}=-\sqrt{3}\cdot\dfrac{x_1+x_2-2x_0}{x_1-x_2}$，直线$PM: y=-\sqrt{3}(x-x_0)+y_0$，即$y=y_0+\sqrt{3}x_0-\sqrt{3}x$，代入双曲线的方程$3x^2-y^2-3=0$，即$(\sqrt{3}x+y)(\sqrt{3}x-y)=3$中，得：$(y_0+\sqrt{3}x_0)[2\sqrt{3}x-(y_0+\sqrt{3}x_0)]=3$，解得$P$的横坐标$x_1=\dfrac{1}{2\sqrt{3}}\left(\dfrac{3}{y_0+\sqrt{3}x_0}+y_0+\sqrt{3}x_0\right)$，同理：$x_2=-\dfrac{1}{2\sqrt{3}}\left(\dfrac{3}{y_0-\sqrt{3}x_0}+y_0-\sqrt{3}x_0\right)$，∴$x_1-x_2=\dfrac{1}{\sqrt{3}}\left(\dfrac{3y_0}{y_0^2-3x_0^2}+y_0\right)$，$x_1+x_2-2x_0=-\dfrac{3x_0}{y_0^2-3x_0^2}-x_0$，∴$m=\dfrac{3x_0}{y_0}$，∴条件②$PQ\parallel AB$等价于$m=k\Leftrightarrow ky_0=3x_0$，综上所述：条件①$M$在$AB$上，等价于$ky_0=k^2(x_0-2)$；条件②$PQ\parallel AB$等价于$ky_0=3x_0$；条件③$|AM|=|BM|$等价于$x_0+ky_0=\dfrac{8k^2}{k^2-3}$．选①②推③：由①②解得$x_0=\dfrac{2k^2}{k^2-3}$，∴$x_0+ky_0=4x_0=\dfrac{8k^2}{k^2-3}$，∴③成立；选①③推②：由①③解得$x_0=\dfrac{2k^2}{k^2-3}$，$ky_0=\dfrac{6k^2}{k^2-3}$，∴$ky_0=3x_0$，∴②成立；选②③推①：由②③解得$x_0=\dfrac{2k^2}{k^2-3}$，$ky_0=\dfrac{6k^2}{k^2-3}$，∴$x_0-2=\dfrac{6}{k^2-3}$，∴$ky_0=k^2(x_0-2)$，∴①成立.
\end{solution}
\end{question}

\begin{question}[points=12]
已知函数$f(x)=xe^{ax}-e^x$．

(1)当$a=1$时，讨论$f(x)$的单调性；

(2)当$x>0$时，$f(x)<-1$，求$a$的取值范围；

(3)设$n\in\mathbf{N}^*$，证明：$\dfrac{1}{1^2+1}+\dfrac{1}{2^2+2}+\cdots+\dfrac{1}{n^2+n}>\ln(n+1)$．
\begin{solution}
\textbf{答案：}（1）$f(x)$的减区间为$(-\infty,0)$，增区间为$(0,+\infty)$；（2）$a\leq\dfrac{1}{2}$；（3）见解析

\textbf{分析：}（1）求出$f'(x)$，讨论其符号后可得$f(x)$的单调性.（2）设$h(x)=xe^{ax}-e^x+1$，求出$h''(x)$，先讨论$a>\dfrac{1}{2}$时题设中的不等式不成立，再就$0<a\leq\dfrac{1}{2}$结合放缩法讨论$h'(x)$符号，最后就$a\leq 0$结合放缩法讨论$h(x)$的范围后可得参数的取值范围.（3）由（2）可得$2\ln t<t-\dfrac{1}{t}$对任意的$t>1$恒成立，从而可得$\ln(n+1)-\ln n<\dfrac{1}{n^2+n}$对任意的$n\in\mathbf{N}^*$恒成立，结合裂项相消法可证题设中的不等式.

\textbf{详解：}（1）当$a=1$时，$f(x)=(x-1)e^x$，则$f'(x)=xe^x$，当$x<0$时，$f'(x)<0$，当$x>0$时，$f'(x)>0$，故$f(x)$的减区间为$(-\infty,0)$，增区间为$(0,+\infty)$.
（2）设$h(x)=xe^{ax}-e^x+1$，则$h(0)=0$，又$h'(x)=(1+ax)e^{ax}-e^x$，设$g(x)=(1+ax)e^{ax}-e^x$，则$g'(x)=(2a+a^2x)e^{ax}-e^x$，若$a>\dfrac{1}{2}$，则$g'(0)=2a-1>0$，因为$g'(x)$为连续不间断函数，故存在$x_0>0$，使得$\forall x\in(0,x_0)$，总有$g'(x)>0$，故$g(x)$在$(0,x_0)$为增函数，故$g(x)>g(0)=0$，故$h(x)$在$(0,x_0)$为增函数，故$h(x)>h(0)=0$，即$f(x)>-1$，与题设矛盾. 若$0<a\leq\dfrac{1}{2}$，则$h'(x)=(1+ax)e^{ax}-e^x=e^{ax+\ln(1+ax)}-e^x$，下证：对任意$x>0$，总有$\ln(1+x)<x$成立，证明：设$S(x)=\ln(1+x)-x$，故$S'(x)=\dfrac{1}{1+x}-1=\dfrac{-x}{1+x}<0$，故$S(x)$在$(0,+\infty)$上为减函数，故$S(x)<S(0)=0$即$\ln(1+x)<x$成立. 由上述不等式有$e^{ax+\ln(1+ax)}-e^x<e^{ax+ax}-e^x=e^{2ax}-e^x\leq 0$（因为$2a\leq 1$，且$x>0$时$e^{2ax}\leq e^x$），故$h'(x)\leq 0$总成立，即$h(x)$在$(0,+\infty)$上为减函数，所以$h(x)<h(0)=0$，即$f(x)<-1$. 当$a\leq 0$时，有$h'(x)=e^{ax}-e^x+axe^{ax}<1-1+0=0$（因为$e^{ax}\leq 1$，$e^x>1$，$axe^{ax}\leq 0$），所以$h(x)$在$(0,+\infty)$上为减函数，所以$h(x)<h(0)=0$，即$f(x)<-1$. 综上，$a\leq\dfrac{1}{2}$.
（3）取$a=\dfrac{1}{2}$，则$\forall x>0$，总有$xe^{\frac{x}{2}}-e^x+1<0$成立，令$t=e^{\frac{x}{2}}$，则$t>1$，$t^2=e^x$，$x=2\ln t$，故$2t\ln t<t^2-1$即$2\ln t<t-\dfrac{1}{t}$对任意的$t>1$恒成立. 所以对任意的$n\in\mathbf{N}^*$，有$2\ln\dfrac{n+1}{n}<\dfrac{n+1}{n}-\dfrac{n}{n+1}$，整理得到$\ln(n+1)-\ln n<\dfrac{1}{n^2+n}$，故$\dfrac{1}{1^2+1}+\dfrac{1}{2^2+2}+\cdots+\dfrac{1}{n^2+n}>\ln2-\ln1+\ln3-\ln2+\cdots+\ln(n+1)-\ln n=\ln(n+1)$，故不等式成立.
\end{solution}
\end{question}

